\documentclass[12pt]{article}
\usepackage[T1]{fontenc}
\usepackage[utf8]{inputenc}
\usepackage{lmodern}
\usepackage{textcomp}
\usepackage{enumitem}
\usepackage{geometry}
\geometry{margin=1in}
\begin{document}
\begin{center}
Answer Key - Psychology - Undergraduate Level Assessment 18 \
\small (Answers are ordered exactly as in the question paper.)
\end{center}

\section*{Multiple Choice}
\begin{enumerate}[label=\arabic*.]
\item \textbf{Answer:} A
\\ \textit{Options:} A) Onset phase, abuse phase, dependency phase, abstinence phase; B) Prealcoholic phase, alcoholic phase, recovery phase, relapse phase; C) Initial phase, adaptive phase, dependent phase, deteriorative phase; D) Experimental phase, social use phase, intensive use phase, compulsive use phase; E) Initial phase, middle phase, crucial phase, termination phase
\\ \textit{Note:} Correct option: A - Onset phase, abuse phase, dependency phase, abstinence phase
\item \textbf{Answer:} E
\\ \textit{Options:} A) Arithmetic and Digit Span subtests; B) All subtests except Block Design; C) Similarities and Picture Completion subtests; D) Vocabulary subtest only; E) All Performance subtests
\\ \textit{Note:} Correct option: E - All Performance subtests
\item \textbf{Answer:} B
\\ \textit{Options:} A) Motivation can exist without any form of arousal.; B) Arousal level is a central aspect of motivation.; C) Motivation is the cause of arousal.; D) Motivation is a by-product of arousal.; E) High arousal always maximizes motivation.
\\ \textit{Note:} Correct option: B - Arousal level is a central aspect of motivation.
\item \textbf{Answer:} A
\\ \textit{Options:} A) decreasing both within-group and between-group variability.; B) increasing both within-group and between-group variability.; C) increasing between-group variability only.; D) increasing within-group variability.; E) decreasing between-group variability only.
\\ \textit{Note:} Correct option: A - decreasing both within-group and between-group variability.
\item \textbf{Answer:} B
\\ \textit{Options:} A) "the message is within his "latitude of acceptance."""; B) he is in a neutral or slightly negative mood.; C) the message is outside his "latitude of acceptance."; D) everyone else in the group disagrees with the message.; E) everyone else in the group agrees with the message.
\\ \textit{Note:} Correct option: B - he is in a neutral or slightly negative mood.
\end{enumerate}

\section*{True / False}
\begin{enumerate}[label=\arabic*.]
\setcounter{enumi}{5}
\item \textbf{Answer:} False
\\ \textit{Options:} A) True; B) False
\\ \textit{Note:} LLM-generated false statement. Correct answer: 'influencing patients by controlling the consequences of their actions'.
\item \textbf{Answer:} True
\\ \textit{Options:} A) True; B) False
\\ \textit{Note:} LLM-generated true statement based on MCQ.
\item \textbf{Answer:} False
\\ \textit{Options:} A) True; B) False
\\ \textit{Note:} LLM-generated false statement. Correct answer: 'identity vs. role confusion'.
\item \textbf{Answer:} True
\\ \textit{Options:} A) True; B) False
\\ \textit{Note:} LLM-generated true statement based on MCQ.
\item \textbf{Answer:} True
\\ \textit{Options:} A) True; B) False
\\ \textit{Note:} LLM-generated true statement based on MCQ.
\end{enumerate}

\section*{Long Form Response}
\begin{enumerate}[label=\arabic*.]
\setcounter{enumi}{10}
\item \textbf{Answer:} ask more questions.. Explain why this is the correct answer and provide supporting evidence. Reference key facts or concepts that support this choice.
\\ \textit{Note:} Long-form derived from MCQ; award credit for stating the correct idea and giving supporting reasoning.
\item \textbf{Answer:} McClelland's Achievement Motivation model was influenced by Murray's theory of motivation that human behavior is a reflection of individual needs and motives, and that anything leading to positive affect produces approach behavior and anything leading to negative affect produces avoidance behavior.. Explain why this is the correct answer and provide supporting evidence. Reference key facts or concepts that support this choice.
\\ \textit{Note:} Long-form derived from MCQ; award credit for stating the correct idea and giving supporting reasoning.
\end{enumerate}

\vfill
\noindent\textit{End of Answer Key}
\end{document}